\documentclass[a4paper,11pt,french]{article}
\usepackage[utf8]{inputenc}

\usepackage[T1]{fontenc}
\usepackage[francais]{babel} 
\usepackage[top=2cm, bottom=2cm, left=2cm, right=2cm, includeheadfoot]{geometry} %pour les marges
\usepackage{lmodern}
\usepackage{pict2e}
\usepackage{tikz}	
\usepackage{tikz-uml}
\usepackage{fancyhdr} % Required for custom headers
\usepackage{lastpage} % Required to determine the last page for the footer
\usepackage{extramarks} % Required for headers and footers
\usepackage{graphicx} % Required to insert images
\usepackage{tabularx, longtable}
\usepackage{color, colortbl}
\usepackage{lscape}
%\usepackage[hidelinks]{hyperref}
\usepackage{longtable}
\usepackage{multirow}
\usepackage{rotating}
%\usepackage{pgfgantt}
%\usepackage{pgfcalendar}
%\usepackage{ifthen}
\usepackage{gensymb}
\usepgflibrary{arrows} % for pgf-umlsd

\usetikzlibrary{trees,shapes.geometric,arrows,decorations.pathmorphing,backgrounds,fit,positioning,shapes.symbols,chains	}

\linespread{1.1} % Line spacing

% Set up the header and footer
\pagestyle{fancy}
\lhead{\textbf{\hmwkClass -- \hmwkSubject \\ \hmwkTitle \\ \hmwkDocName}} % Top left header
\rhead{\includegraphics[width=10em]{logo_univ.png}}
\lfoot{\lastxmark} % Bottom left footer
\cfoot{} % Bottom center footer
\rfoot{Page\ \thepage\ / \pageref{LastPage}} % Bottom right footer
\renewcommand\headrulewidth{0.4pt} % Size of the header rule
\renewcommand\footrulewidth{0.4pt} % Size of the footer rule

\setlength{\headheight}{40pt}

\newcommand{\hmwkTitle}{Transchiffrement} % Assignment title
\newcommand{\hmwkClass}{Master 2 SSI } % Course/class
\newcommand{\hmwkAuthorName}{Emile Générat} % Your name
\newcommand{\hmwkSubject}{Conduite de projet} % Subject
\newcommand{\hmwkDocName}{Analyse des Risques} % Document name

\newcommand{\version}{1.0} % Document version
\newcommand{\docDate}{13 Décembre 2013} % Document date
\newcommand{\checked}{} % Checker name
\newcommand{\approved}{} % Approver name

\makeatletter
\newcommand{\resettranslate}{\let\translate\@firstofone}
\makeatother

\definecolor{gris}{rgb}{0.95, 0.95, 0.95}

\title{
\vspace{2in}
\textmd{\textbf{\hmwkClass :\ \hmwkTitle}}\\
\normalsize\vspace{0.1in}\small{Due\ on\ \hmwkDueDate}\\
\vspace{0.1in}\large{\textit{\hmwkClassInstructor\ \hmwkClassTime}}
\vspace{3in}
}

\author{\hmwkAuthorName}
\date{} % Insert date here if you want it to appear below your name


\usepackage{amsmath}
\begin{document}
\newcount\startdate
\newcount\daynum
%\pgfcalendardatetojulian{2013-01-021}{\startdate}
\pagestyle{fancy}

\vspace*{5cm}
\begin{center}\textbf{\Huge{\hmwkDocName}}\end{center}
\vspace*{4.5cm}
	

\fcolorbox{black}{gris}{
\begin{minipage}{15cm}
\begin{tabularx}{10cm}{lXl}
	\bfseries{Version} & & \version\\
	& & \\
	\bfseries{Date} & & \docDate\\
	& & \\
	\bfseries{Rédigé par} & & \hmwkAuthorName \\
	& & \\
	\bfseries{Relu par} & & \checked \\
	& & \\
	\bfseries{Approuvé par} & & \approved \\
	& & \\
\end{tabularx}
\end{minipage}
}

\newpage

%Tableau de mises à jour
\vspace*{1cm}
\begin{center}
\textbf{\huge{Versions}}\\
\vspace*{3cm}
	\begin{tabularx}{16cm}{|c|c|X|}
	\hline
	\bfseries{Version} & \bfseries{Date} & \bfseries{Modifications réalisées}\\
	\hline
	1.0 & 13/12/2013 & Création\\
	\hline
	\end{tabularx}
\end{center}

%La table des matières
\clearpage
\tableofcontents
\clearpage


\newpage
\section{Documents de référence}
\begin{itemize}
\item Plan de Développement v1.0
\end{itemize}


\newpage
\section{Portefeuille de risques}

\begin{landscape}

\begin{longtable}{|m{1.5cm}|m{4cm}|m{3cm}|m{3.2cm}|m{2.2cm}|m{2cm}|m{2cm}|m{2.2cm}|}
\hline
\rowcolor[gray]{.8}
Numéro & Risque & Date de création & Pilote & Catégorie \footnotemark[1] & Indice de probabilité & Indice de gravité & Criticité \footnotemark[2] \\ 
\hline
\endfirsthead
\hline
\rowcolor[gray]{.8}
Numéro & Risque & Date de création & Pilote & Catégorie \footnotemark[1] & Indice de probabilité & Indice d'impact & Criticité \footnotemark[2] \\ 
\hline

\endhead
n01 & Ne pas livrer le client à temps & 13/12/2013 &  & Planning & 1 & 3 & A surveiller \\
\hline
n02 & Insatisfaction du client & 13/12/2013 & & Planning & 1 & 4 &  Critique \\
\hline
n03 & Retard, algorithme prêt trop tard & 13/12/2013 & & Planning & 1 & 3 & A surveiller \\
\hline
n04 & Niveau technique insuffisant & 13/12/2013 &  &  Compétence & 1 & 2 & Acceptable \\
\hline
n05 & Implication insuffisante, ou absences & 13/12/2013 &  & Planning & 1 & 3 & A surveiller \\
\hline
n06 & Perte des données du projet & 13/12/2013 &  & Technique & 1 & 4 & Critique \\
\hline
n07 & Performances insuffisantes (transchiffrement trop long) & 13/12/2013 &  & Technique & 1 & 3 & A surveiller \\
\hline
n08 & Modification du Cahier des Charges & 13/12/2013 &  & Technique & 2 & 2 & Acceptable \\
\hline

\end{longtable}
\footnotetext[1]{Compétences / Planning / Technique}
\footnotetext[2]{Acceptable / À surveiller / Critique}

\end{landscape}

\newpage




\begin{center}
\begin{tabular}{|>{\columncolor[gray]{.8}}m{8cm}|m{8cm}|}
\hline
 Intitulé du risque & Ne pas livrer le client à temps \\
\hline
 Description du risque & Un dérapage sur le planning ne nous permet pas de rendre tous les livrables au moment voulu  \\
\hline
Pilote du risque &  \\
\hline
\end{tabular}
\end{center}

\begin{center}
\begin{tabular}{|>{\columncolor[gray]{.8}}m{3.8cm}|m{3.8cm}|>{\columncolor[gray]{.8}}m{3.8cm}|m{3.8cm}|}
\hline
Indice de gravité & 3 &Date de début d'effet& Début de la conception \\
\hline
Indice de probabilité & 1 & Date de fin d'effet & Fin du projet\\
\hline
Criticité \footnotemark[1] & À surveiller & Découvreur & \\
\hline
État \footnotemark[2] & Actif & Catégorie \footnotemark[3] & Planning\\
\hline
\end{tabular}
\end{center}

\begin{center}
\begin{tabular}{|m{5cm}|m{11cm}|}
\hline
\rowcolor[gray]{.8} Actions préventives & Description\\
\hline
 Réunions quotidiennes & Pour ne pas déborder du planning, cette réunion nous permettra de ne pas nous éparpiller \\
\hline
\end{tabular}
\end{center}

\begin{center}
\begin{tabular}{|m{5cm}|m{11cm}|}
\hline
\rowcolor[gray]{.8} Actions correctives & Description\\
\hline
%\aRemplir{Nom de l'action} & \aRemplir{Lister les actions corrective en étant le plus précis possible} \\
Abandonner une fonctionnalité optionnelle & Pour ne pas dépasser le temps imparti, une fonctionnalité optionnelle non-finalisée nous permettra de nous focaliser sur le sujet principal \\
\hline
\end{tabular}
\end{center}

\begin{center}
\begin{tabular}{|m{2.6cm}|m{5cm}|m{3.8cm}|m{3.8cm}|}
\hline
\multicolumn{4}{|>{\columncolor[gray]{.8}}c|}{Suivi du risque} \\
\hline
\hline
\rowcolor[gray]{.8} Date & Description & Signature du pilote de risque & Signature du Chef de Projet \\
\hline
 &  &  & \\
\hline
\end{tabular}
\end{center}

\footnotetext[1]{Acceptable / À surveiller / Critique}
\footnotetext[2]{Actif / Inactif}
\footnotetext[3]{Compétences / Planning / Technique}

\newpage



\begin{center}
\begin{tabular}{|>{\columncolor[gray]{.8}}m{8cm}|m{8cm}|}
\hline
 Intitulé du risque &  Insatisfaction du client\\
\hline
 Description du risque & Malgré la finition du projet, le rendu n'est pas ce à quoi s'attendait le client  \\
\hline
Pilote du risque &  \\
\hline
\end{tabular}
\end{center}

\begin{center}
\begin{tabular}{|>{\columncolor[gray]{.8}}m{3.8cm}|m{3.8cm}|>{\columncolor[gray]{.8}}m{3.8cm}|m{3.8cm}|}
\hline
Indice de gravité & 1 &Date de début d'effet& Début de la conception \\
\hline
Indice de probabilité & 4 & Date de fin d'effet & Fin du projet\\
\hline
Criticité \footnotemark[1] & Critique & Découvreur & \\
\hline
État \footnotemark[2] & Actif & Catégorie \footnotemark[3] & Planning\\
\hline
\end{tabular}
\end{center}

\begin{center}
\begin{tabular}{|m{5cm}|m{11cm}|}
\hline
\rowcolor[gray]{.8} Actions préventives & Description\\
\hline
Réunions hebdomadaires avec le client & Pour éviter de partir dans la mauvaise direction, ces réunions avec le client nous aideront à respecter ses attentes \\
\hline
\end{tabular}
\end{center}

\begin{center}
\begin{tabular}{|m{5cm}|m{11cm}|}
\hline
\rowcolor[gray]{.8} Actions correctives & Description\\
\hline
%\aRemplir{Nom de l'action} & \aRemplir{Lister les actions corrective en étant le plus précis possible} \\
&\\
\hline
\end{tabular}
\end{center}

\begin{center}
\begin{tabular}{|m{2.6cm}|m{5cm}|m{3.8cm}|m{3.8cm}|}
\hline
\multicolumn{4}{|>{\columncolor[gray]{.8}}c|}{Suivi du risque} \\
\hline
\hline
\rowcolor[gray]{.8} Date & Description & Signature du pilote de risque & Signature du Chef de Projet \\
\hline
 &  &  & \\
\hline
\end{tabular}
\end{center}

\footnotetext[1]{Acceptable / À surveiller / Critique}
\footnotetext[2]{Actif / Inactif}
\footnotetext[3]{Compétences / Planning / Technique}

\newpage

\begin{center}
\begin{tabular}{|>{\columncolor[gray]{.8}}m{8cm}|m{8cm}|}
\hline
 Intitulé du risque &  Retard, algorithme prêt trop tard\\
\hline
 Description du risque & L'algorithme de recherche n'a pas été fini assez tôt à cause de sa complexité, ce qui n'a pas pu donner lieu à un résultat que nous aurions voulu utiliser.  \\
\hline
Pilote du risque &  \\
\hline
\end{tabular}
\end{center}

\begin{center}
\begin{tabular}{|>{\columncolor[gray]{.8}}m{3.8cm}|m{3.8cm}|>{\columncolor[gray]{.8}}m{3.8cm}|m{3.8cm}|}
\hline
Indice de gravité & 1 &Date de début d'effet& Début de la conception \\
\hline
Indice de probabilité & 3 & Date de fin d'effet & Fin du projet\\
\hline
Criticité \footnotemark[1] & A surveiller & Découvreur & \\
\hline
État \footnotemark[2] & Actif & Catégorie \footnotemark[3] & Planning\\
\hline
\end{tabular}
\end{center}

\begin{center}
\begin{tabular}{|m{5cm}|m{11cm}|}
\hline
\rowcolor[gray]{.8} Actions préventives & Description\\
\hline
Réunions hebdomadaires avec le client & Pour éviter de partir dans la mauvaise direction, ces réunions avec le client nous aideront à respecter ses attentes \\
\hline
Recherche & Travail de compréhension de l'attaque sur MD5 à partir d'articles publiés et récupération d'aide sur l'algorithme utilisé dans cette recherche. \\
\hline
\end{tabular}
\end{center}

\begin{center}
\begin{tabular}{|m{5cm}|m{11cm}|}
\hline
\rowcolor[gray]{.8} Actions correctives & Description\\
\hline
%\aRemplir{Nom de l'action} & \aRemplir{Lister les actions corrective en étant le plus précis possible} \\
& \\
\hline
\end{tabular}
\end{center}

\begin{center}
\begin{tabular}{|m{2.6cm}|m{5cm}|m{3.8cm}|m{3.8cm}|}
\hline
\multicolumn{4}{|>{\columncolor[gray]{.8}}c|}{Suivi du risque} \\
\hline
\hline
\rowcolor[gray]{.8} Date & Description & Signature du pilote de risque & Signature du Chef de Projet \\
\hline
 &  &  & \\
\hline
\end{tabular}
\end{center}


\footnotetext[1]{Acceptable / À surveiller / Critique}
\footnotetext[2]{Actif / Inactif}
\footnotetext[3]{Compétences / Planning / Technique}

\newpage

\begin{center}
\begin{tabular}{|>{\columncolor[gray]{.8}}m{8cm}|m{8cm}|}
\hline
 Intitulé du risque &  Niveau technique insuffisant\\
\hline
 Description du risque & Un manque de technique d'un ou plusieurs des étudiants nous handicape et ne nous permet pas de réaliser le projet en totalité  \\
\hline
Pilote du risque &  \\
\hline
\end{tabular}
\end{center}

\begin{center}
\begin{tabular}{|>{\columncolor[gray]{.8}}m{3.8cm}|m{3.8cm}|>{\columncolor[gray]{.8}}m{3.8cm}|m{3.8cm}|}
\hline
Indice de gravité & 1 &Date de début d'effet& Début de la conception \\
\hline
Indice de probabilité & 2 & Date de fin d'effet & Fin du projet\\
\hline
Criticité \footnotemark[1] & Acceptable & Découvreur & \\
\hline
État \footnotemark[2] & Actif & Catégorie \footnotemark[3] & Compétence\\
\hline
\end{tabular}
\end{center}

\begin{center}
\begin{tabular}{|m{5cm}|m{11cm}|}
\hline
\rowcolor[gray]{.8} Actions préventives & Description\\
\hline
Apprentissage & Que ce soit par le biais de tutoriels ou tout simplement d'une explication d'un autre étudiant, l'apprentissage et l'approfondissement des connaissances devra être fait en amont.\\
\hline
\end{tabular}
\end{center}

\begin{center}
\begin{tabular}{|m{5cm}|m{11cm}|}
\hline
\rowcolor[gray]{.8} Actions correctives & Description\\
\hline
%\aRemplir{Nom de l'action} & \aRemplir{Lister les actions corrective en étant le plus précis possible} \\
Entre-aide & Si un étudiant n'est pas capable de développer sa partie à cause d'un manque de connaissance, il sera aidé par un autre étudiant sur l'ensemble de sa tâche ce qui lui fera comprendre comment il faut faire. \\
\hline
\end{tabular}
\end{center}

\begin{center}
\begin{tabular}{|m{2.6cm}|m{5cm}|m{3.8cm}|m{3.8cm}|}
\hline
\multicolumn{4}{|>{\columncolor[gray]{.8}}c|}{Suivi du risque} \\
\hline
\hline
\rowcolor[gray]{.8} Date & Description & Signature du pilote de risque & Signature du Chef de Projet \\
\hline
 &  &  & \\
\hline
\end{tabular}
\end{center}


\footnotetext[1]{Acceptable / À surveiller / Critique}
\footnotetext[2]{Actif / Inactif}
\footnotetext[3]{Compétences / Planning / Technique}

\newpage


\begin{center}
\begin{tabular}{|>{\columncolor[gray]{.8}}m{8cm}|m{8cm}|}
\hline
 Intitulé du risque &  Implication insuffisante ou absences \\
\hline
 Description du risque & Un ou plusieurs étudiants n'accomplissent pas les tâches qui leur ont été données ou sont absents trop longtemps pour ne pas avoir le temps de les faire. \\
\hline
Pilote du risque &  \\
\hline
\end{tabular}
\end{center}

\begin{center}
\begin{tabular}{|>{\columncolor[gray]{.8}}m{3.8cm}|m{3.8cm}|>{\columncolor[gray]{.8}}m{3.8cm}|m{3.8cm}|}
\hline
Indice de gravité & 1 &Date de début d'effet& Début de la conception \\
\hline
Indice de probabilité & 3 & Date de fin d'effet & Fin du projet\\
\hline
Criticité \footnotemark[1] & A surveiller & Découvreur & \\
\hline
État \footnotemark[2] & Actif & Catégorie \footnotemark[3] & Planning\\
\hline
\end{tabular}
\end{center}

\begin{center}
\begin{tabular}{|m{5cm}|m{11cm}|}
\hline
\rowcolor[gray]{.8} Actions préventives & Description\\
\hline
Réunions quotidiennes & Pour surveiller que tout le monde accomplit bien les tâches qu'on leur donne, une réunion quotidienne permettra de sentir la motivation des étudiants.\\
\hline
Prévoir les absences certaines & Pour que les absences ne soient pas dérangeantes pour le projet (absences des deux personnes qui travaillent sur le même point), les absences seront prévues et annoncées à l'avance\\
\hline
\end{tabular}
\end{center}

\begin{center}
\begin{tabular}{|m{5cm}|m{11cm}|}
\hline
\rowcolor[gray]{.8} Actions correctives & Description\\
\hline
%\aRemplir{Nom de l'action} & \aRemplir{Lister les actions corrective en étant le plus précis possible} \\
Réaffectation des tâches & Un étudiant qui peine à travailler sur une tâche car elle ne l'intéresse pas pourra se voir affecter à une autre tâche pour tenter de le remotiver.\\
\hline
Sanction & Un étudiant qui est trop absent sera amené à s'expliquer avec le professeur encadrant et sera bien sûr réprimandé.\\
\hline
\end{tabular}
\end{center}

\begin{center}
\begin{tabular}{|m{2.6cm}|m{5cm}|m{3.8cm}|m{3.8cm}|}
\hline
\multicolumn{4}{|>{\columncolor[gray]{.8}}c|}{Suivi du risque} \\
\hline
\hline
\rowcolor[gray]{.8} Date & Description & Signature du pilote de risque & Signature du Chef de Projet \\
\hline
 &  &  & \\
\hline
\end{tabular}
\end{center}


\footnotetext[1]{Acceptable / À surveiller / Critique}
\footnotetext[2]{Actif / Inactif}
\footnotetext[3]{Compétences / Planning / Technique}

\newpage


\begin{center}
\begin{tabular}{|>{\columncolor[gray]{.8}}m{8cm}|m{8cm}|}
\hline
 Intitulé du risque &  Perte des données du projet \\
\hline
 Description du risque & Une mauvaise manipulation ou un crash du serveur où sont stockées les données résultent en une perte de toute les données. \\
\hline
Pilote du risque &  \\
\hline
\end{tabular}
\end{center}

\begin{center}
\begin{tabular}{|>{\columncolor[gray]{.8}}m{3.8cm}|m{3.8cm}|>{\columncolor[gray]{.8}}m{3.8cm}|m{3.8cm}|}
\hline
Indice de gravité & 1 &Date de début d'effet& Début de la conception \\
\hline
Indice de probabilité & 4 & Date de fin d'effet & Fin du projet\\
\hline
Criticité \footnotemark[1] & Critique & Découvreur & \\
\hline
État \footnotemark[2] & Actif & Catégorie \footnotemark[3] & Technique\\
\hline
\end{tabular}
\end{center}

\begin{center}
\begin{tabular}{|m{5cm}|m{11cm}|}
\hline
\rowcolor[gray]{.8} Actions préventives & Description\\
\hline
Sauvegarde & Tous les soirs, une copie des données sera faite en local sur chaque ordinateur, ce qui permettra de recréer le projet à partir de ces données \\
\hline
\end{tabular}
\end{center}

\begin{center}
\begin{tabular}{|m{5cm}|m{11cm}|}
\hline
\rowcolor[gray]{.8} Actions correctives & Description\\
\hline
%\aRemplir{Nom de l'action} & \aRemplir{Lister les actions corrective en étant le plus précis possible} \\
Récupération des données & Utilisation de la sauvegarde locale sur un des ordinateurs du groupe \\
\hline
\end{tabular}
\end{center}

\begin{center}
\begin{tabular}{|m{2.6cm}|m{5cm}|m{3.8cm}|m{3.8cm}|}
\hline
\multicolumn{4}{|>{\columncolor[gray]{.8}}c|}{Suivi du risque} \\
\hline
\hline
\rowcolor[gray]{.8} Date & Description & Signature du pilote de risque & Signature du Chef de Projet \\
\hline
 &  &  & \\
\hline
\end{tabular}
\end{center}


\footnotetext[1]{Acceptable / À surveiller / Critique}
\footnotetext[2]{Actif / Inactif}
\footnotetext[3]{Compétences / Planning / Technique}

\newpage


\begin{center}
\begin{tabular}{|>{\columncolor[gray]{.8}}m{8cm}|m{8cm}|}
\hline
 Intitulé du risque &  Performances insuffisantes \\
\hline
 Description du risque & Une mauvaise implémentation du code ne satisfait pas les contraintes de rapidité pour rendre le transchiffrement du proxy transparent. \\
\hline
Pilote du risque &  \\
\hline
\end{tabular}
\end{center}

\begin{center}
\begin{tabular}{|>{\columncolor[gray]{.8}}m{3.8cm}|m{3.8cm}|>{\columncolor[gray]{.8}}m{3.8cm}|m{3.8cm}|}
\hline
Indice de gravité & 1 &Date de début d'effet& Début de la conception \\
\hline
Indice de probabilité & 3 & Date de fin d'effet & Fin du projet\\
\hline
Criticité \footnotemark[1] & A surveiller & Découvreur & \\
\hline
État \footnotemark[2] & Actif & Catégorie \footnotemark[3] & Technique\\
\hline
\end{tabular}
\end{center}

\begin{center}
\begin{tabular}{|m{5cm}|m{11cm}|}
\hline
\rowcolor[gray]{.8} Actions préventives & Description\\
\hline
Tests d'efficacité & Plusieurs tests seront faits pour que les temps ne dépassent pas ceux voulus. \\
\hline
\end{tabular}
\end{center}

\begin{center}
\begin{tabular}{|m{5cm}|m{11cm}|}
\hline
\rowcolor[gray]{.8} Actions correctives & Description\\
\hline
%\aRemplir{Nom de l'action} & \aRemplir{Lister les actions corrective en étant le plus précis possible} \\
Optimisation des algorithmes & Puisqu'un gain de temps est nécessaire, la révision des algorithmes développés et une optimisation de ces codes seront faites.  \\
\hline
\end{tabular}
\end{center}

\begin{center}
\begin{tabular}{|m{2.6cm}|m{5cm}|m{3.8cm}|m{3.8cm}|}
\hline
\multicolumn{4}{|>{\columncolor[gray]{.8}}c|}{Suivi du risque} \\
\hline
\hline
\rowcolor[gray]{.8} Date & Description & Signature du pilote de risque & Signature du Chef de Projet \\
\hline
 &  &  & \\
\hline
\end{tabular}
\end{center}


\footnotetext[1]{Acceptable / À surveiller / Critique}
\footnotetext[2]{Actif / Inactif}
\footnotetext[3]{Compétences / Planning / Technique}

\newpage

\begin{center}
\begin{tabular}{|>{\columncolor[gray]{.8}}m{8cm}|m{8cm}|}
\hline
 Intitulé du risque &  Modification du cachier des charges \\
\hline
 Description du risque & Le client souhaite modifier ou ajouter des fonctionnalités au cahier des charges  \\
\hline
Pilote du risque &  \\
\hline
\end{tabular}
\end{center}

\begin{center}
\begin{tabular}{|>{\columncolor[gray]{.8}}m{3.8cm}|m{3.8cm}|>{\columncolor[gray]{.8}}m{3.8cm}|m{3.8cm}|}
\hline
Indice de gravité & 2 &Date de début d'effet& Début de la conception \\
\hline
Indice de probabilité & 2 & Date de fin d'effet & Fin du projet\\
\hline
Criticité \footnotemark[1] & Acceptable & Découvreur & \\
\hline
État \footnotemark[2] & Actif & Catégorie \footnotemark[3] & Technique\\
\hline
\end{tabular}
\end{center}

\begin{center}
\begin{tabular}{|m{5cm}|m{11cm}|}
\hline
\rowcolor[gray]{.8} Actions préventives & Description\\
\hline
 Réunions hebdomadaires avec le client & Pour éviter de créer de nouveaux besoins, ces réunions avec le client nous aideront à respecter ses attentes et à limiter les modifications \\
\hline
\end{tabular}
\end{center}

\begin{center}
\begin{tabular}{|m{5cm}|m{11cm}|}
\hline
\rowcolor[gray]{.8} Actions correctives & Description\\
\hline
%\aRemplir{Nom de l'action} & \aRemplir{Lister les actions corrective en étant le plus précis possible} \\
Prise en compte des modifications & Une réorganisation du planning sera faite si nous avons le temps et les moyens de réaliser cette modification, sinon cette modification sera considérée d'un commun accord avec le client, comme optionnelle\\
\hline
\end{tabular}
\end{center}

\begin{center}
\begin{tabular}{|m{2.6cm}|m{5cm}|m{3.8cm}|m{3.8cm}|}
\hline
\multicolumn{4}{|>{\columncolor[gray]{.8}}c|}{Suivi du risque} \\
\hline
\hline
\rowcolor[gray]{.8} Date & Description & Signature du pilote de risque & Signature du Chef de Projet \\
\hline
 &  &  & \\
\hline
\end{tabular}
\end{center}


\footnotetext[1]{Acceptable / À surveiller / Critique}
\footnotetext[2]{Actif / Inactif}
\footnotetext[3]{Compétences / Planning / Technique}

\end{document}